\subsection{Prospective W06: certified opportunity, stopping and multi-day economics}
\label{sec:w06}
\status{Protocol only. No W06 experiment or Lean build is claimed. Freeze decisions and inputs before holdout access; preserve all W05 evidence.}
The new aim is not another compression competition. It is to discover whether a strong, reliable baseline leaves enough opportunity to merit optimization. W05 and the separate retrospective economic diagnostic are development and regression inputs, never W06 confirmatory holdouts. Prioritize B/A jointly (mathematical and numerical certificates), then E (economic control where opportunity remains), then C (transfer to static AC and UC). Retire W06-D grouping unless a genuinely distinct, separately justified mechanism is proposed; PE/HIL remains a separate unrun study.

\begin{longtable}{P{.14\linewidth}P{.43\linewidth}P{.34\linewidth}}
\toprule Track & Design and mandatory comparator & Admission gate / falsifier\\\midrule\endhead
W06-B (critical) & Formalize single-case and weighted-panel COA inequalities, strict/weak stop thresholds, boxed LP dual signs, nesting, net-benefit ceiling and conditional candidate feasibility. Include exact positive/negative controls for invalid relaxation, sign, units, infeasible baseline and missing scenario. & Actual Lean build, pinned dependency and axiom inventory. No theorem implies population or plant assurance. Counterexample to any stronger claim must be preserved.\\
W06-A (critical) & Compare increasingly strong lower bounds on identical fresh service/cost storage instances: elementary energy/ramp bounds, unrestricted LP, tighter relaxations if worthwhile, full-model bounds. Compare strong baseline, solver-native bounds and exact witness checks. & Bound validity and total time/cost; fraction/weight of targets excluded; no false exclusion on known-improvement controls; report unresolved intervals and missed deadlines.\\
W06-E (critical) & Chronological multi-day storage/compute policies under varying prices, reserve/outage, workload service, degradation and inventory. Strong feasible fill/price-aware rule, rolling LP/MILP and proposed controller with common issued information. Offline references separate. & Bound opportunity first; preserve service and risk or explicitly price approved tradeoffs. Only test candidates where material opportunity remains. Positive held-out net benefit with validated reliability and complete clocks; no extrapolation.\\
W06-C (high) & Static PGLib-OPF and PGLib-UC cost opportunity bounds: strong feasible baseline; lossless/convex relaxation where valid, full UC LP and independently valid full-model bound. New cases, pins, multiple regimes, failed solves retained. & Valid lower/upper interval under identical model. No local AC solution labeled globally optimal, no UC mislabeled SCUC, and no unrestricted bound from restricted model. Stop when target excluded.\\\bottomrule
\end{longtable}

\subsubsection{W06-B: theorem and certificate obligations}
State minimization and maximization versions, empty/nonempty domain conditions, feasible incumbent, bounded objectives, attainable versus infimum, and exact numerical semantics. Prove Equation~\ref{eq:coa}, its finite weighted analogue and correct stopping inequality, including equality at the threshold. Prove nesting only under stated inclusion of policy/trajectory sets and common objective; an offline relaxation need not be implementable. For Equation~\ref{eq:coabox}, verify the multiplier signs and every bound endpoint, reconstruct reduced costs exactly, and handle unbounded variables/rows explicitly. Prove a net-benefit exclusion when a known cost floor shares the horizon and units. Supply counterexamples showing why a service ceiling need not bound cost, a positive ceiling need not be attainable, an invalid restriction can create an overstated lower bound, and average savings do not certify tail risk. Attempt a Lean build with pinned inputs, axiom audit and mutation controls, without claiming a new universal theorem. If compilation cannot be executed, report the blocker and retain written derivations only. Review prior art on weak duality, branch-and-bound stopping, value of information, bounds in stochastic programming and value-of-computation; COA's potential contribution is the integrated engineering protocol, not the discovery of duality.

\subsubsection{W06-A: opportunity-screening experiment}
Use W05 regression before freeze, then independent seed families spanning genuine high and low headroom, forced opposition, tight reserve and outages, variable energy price, differing terminal contracts and both feasible and infeasible/missing baselines. Include analytic cases where a known feasible alternative beats the baseline by more than the target, and cases with provable saturation. Register objective units, practical absolute and relative thresholds and eligibility before running holdouts; use paired blocks, keep all cases in the denominator. Validate exact rational lower certificates and exact feasible upper witnesses where available. Where only tolerance-admissible primals exist, label the numerical interval or derive justified feasibility/objective enclosures, not exact $U$. For each tier record $L$, best valid $U$, $U_0-L$, improvement actually witnessed under the admissible policy, target-exclusion coverage, false-exclusion controls, wall time, memory, failure and fallback. An incremental valid lower bound may be retained by max across tiers with the same contract. Compare against a direct strong solver with native proof/bounds and count certificate construction/checking and the cost of finding a valid baseline. Measure marginal coverage per compute cost, including no-gain strata, not only accepted calls. A useful result can be null: direct solving may be cheaper than proving saturation.

\subsubsection{W06-E: inventory carryover and decision value}
First reproduce W05 fill and the two economic diagnostic arms. Keep them historical. For new tests use continuous multi-day horizons (e.g., 7 and 30 days, chosen and justified before holdout), without resetting SOC, outage, ramp history or unresolved workload at midnight. Match day-ahead/rolling issued forecasts; use authorized real traces or fresh synthetic generators clearly distinguished. Register electricity price paths, wear/aging approximation, reserve headroom, charge limits, power and energy outages, protected load, flexible job energy/deadline and service-value assumptions. Do not equate electrical MWh with successful computing jobs without a workload model. Distinguish equal hourly service, equal terminal inventory, minimum stock plus credit, and a genuine service-cost tradeoff as separate contracts. Compare the W05 fill rule, a strong price-aware fill/recharge heuristic, no-continuation rolling optimization, and a development-tuned continuation or robust strategy only after an opportunity certificate leaves material headroom. Retain an H12-information policy that may ignore the tail. Use lexicographic protected-service constraints; any flexible-service tradeoff must be declared and subject to risk noninferiority, headroom and terminal conditions. Offline perfect-foresight LP/MILP bounds are not online performance. Evaluate terminal credit sensitivity or continuation over longer horizons rather than claiming daily drawdown savings; use the full path, record imports, loss, throughput, energy inventory, unmet flexible/protected service, risk, cost, compute, fallback and actuation limitations. Compare out-of-sample chronologically grouped outcomes with paired intervals, practical significance and multiple-comparison controls registered before the run. If the certified opportunity is too small, stop the candidate competition but report the ceiling and its cost. If all policies lose to the strongest rule, preserve the finding.

\subsubsection{W06-C: extend the bound hierarchy to network and commitment}
Use pinned normal/stressed PGLib-OPF cases and UC dates/horizons not used for confirmation previously; match common objective and physical boxes. Baseline feasible AC solutions require independently reconstructed balance and both-end branch/voltage/angle/generator limits. Derive an admissible lossless economic-dispatch floor only where loss/cost monotonicity and other needed assumptions hold; if possible compare established SDP/SOCP or other valid relaxations with verified lower bounds and include validation overhead. Local IPOPT objectives are upper candidates, not lower bounds. If a relaxation is inexact, retain a wide interval rather than claim savings. For UC preserve initial history, start/ramp/reserve, cost and horizon conventions; use strong full MILP and full LP with checked bounds, plus any valid stronger model-specific bounds. Distinguish solver-reported tolerances from exact certificates. A static OPF or global-balance UC model cannot establish SCOPF/SCUC, stability, grid reliability or capacity. Report feasible improvement vs opportunity ceiling separately; do not revisit the failed guided elimination or pairing as a default. Separate case-level upper witnesses, lower certificates and their execution budgets; retain every missing or late result.

\subsubsection{W06 deliverables and stop conditions}
Deliver a standalone W06 technical report and editable TeX, reproduction archive, source locks and hashes, frozen protocol and deviations, actual input/seed and license records, Lean source/logs/axioms if executed, numerical witnesses and bounds, complete per-case/clock/failure records, independent checks, claim-evidence ledger, null findings and blockers. Preserve previous sources and outputs. Do not edit this master or mark any prospective result established. Before claiming any practical benefit, require a valid model boundary, appropriate comparator and information, full validation, risk/terminal equivalence, measured cost and disclosed uncertainty. If the strongest credible bounds exclude a practical gain, that is a successful research result and a reason to retire that narrow objective, not to switch targets after seeing outcomes.
